% 2.4 — Points de restauration, gestion de versions, sauvegardes
\section{Points de restauration, gestion de versions et sauvegardes}
\label{sec:restauration}

\subsection{Les points de restauration des agents}
Claude Code enregistre automatiquement l'état des fichiers modifiés par ses outils
d'édition, et une commande \cmd{/rewind} permet de revenir en arrière. Sa documentation
précise la limite qui intéresse ce projet, avec un exemple sans ambiguïté :
« Checkpointing does not track files modified by Bash commands », illustré par
\cmd{rm file.txt}, \cmd{mv old.txt new.txt} et \cmd{cp source.txt dest.txt}, dont les
modifications « ne peuvent pas être annulées par un retour arrière »~\cite{cc-checkpointing}.
Autrement dit, les commandes shell --- là où se trouvent les destructions --- sont hors du
périmètre.

\subsection{La gestion de versions}
Git protège ce qui a été \emph{validé}. Il ne protège pas les fichiers non suivis ou
ignorés (données locales, fichiers \cmd{.env}, bases SQLite de développement, artefacts),
ni les modifications non validées qu'un \cmd{git reset --hard} ou un \cmd{git checkout --}
efface, ni le dépôt lui-même si l'agent supprime \fichier{.git}. Et un \cmd{git push
--force} réécrit l'historique \emph{du serveur}, que la copie locale ne restaure pas.

\subsection{Les sauvegardes}
Une sauvegarde système (instantané de disque, sauvegarde nocturne) récupère beaucoup, mais
à la mauvaise granularité : elle date de la veille, pas de la seconde qui précède la
commande, et la restaurer est une opération lourde qu'on n'entreprend pas pour un fichier.

\begin{pourquoi}[L'idée qui en découle]
Ce qui manque est une copie \emph{ciblée}, prise \emph{juste avant} la commande, de
\emph{exactement} ce qu'elle va détruire --- et rien qu'elle ne puisse revendiquer à tort.
C'est l'idée de départ de la bêta de \demo{} ; tout le reste du projet consiste à savoir,
avant l'exécution, ce que « exactement ce qu'elle va détruire » veut dire.
\end{pourquoi}
